\documentclass{article}
\usepackage{graphicx} % Required for inserting images
\usepackage{amsmath}
\usepackage{amssymb}

\title{linear-quadratic-regulator}
\author{huangzhengming28@gmail.com}
\date{February 2026}

\begin{document}

\maketitle

\section{Introduction}
This project is the basic implementation of linear-quadratic-regulator(LQR). A rotation object is studied because it is a simple 2-order linear system.

\section{Key Idea}
1. State-Feedback control law;\\
2. Quadratic-Form value function;\\
3. Dynamic-Programming algorithm;\\

\section{Problem Description}
Given a linear time-invariant system:
\begin{equation}
    x_{t+1} = A x_t + B u_t
\end{equation}
and the control Law is:
\begin{equation}
    u_t = -K_t x_t
\end{equation}
to minimize the value function:
\begin{equation}
    J = \sum_{i=0}^{N}{x_t^TQx_t+u_t^TRu_t}
\end{equation}

In summary, the optimization problem is:
\begin{gather}
    \min_{u_0, u_1, ..., u_N}{\sum_{i=0}^{N}{x_t^TQx_t+u_t^TRu_t}} \label{cost_function} \\
    \text{s.t.} \quad x_{t+1} = A x_t + B u_t \label{system_dynamics}\\
    \quad u_t = -K_t x_t \label{control_law}
\end{gather}

\section{Algorithm}

\subsection{Dynamic Programming}
To use Dynamic-Programming algorithm, we first divide the optimization problem into a group of sub problems. For example, for the step $t$:
\begin{gather}
    Q(x_t, u_t) = x_t^T Q x_t + u_t^T R u_t + V(x_{t+1}) \label{sub_problem}\\
    J = \sum_{t=0}^N{Q(x_t, u_t)}
\end{gather}
The algorithm is as follow\\
For $t=N-1, N-2, ..., 0$:\\
1. Solve the optimal $\hat{u_t}$ and transform $Q(x_t, u_t)$ to $V(x_t)$ (which is the optimal value of $\hat{Q(x_t, u_t)}$)\\
2. with $V(x_t)$, $Q(x_{t-1}, u_{t-1})$ can be expressed like \eqref{sub_problem} again

\subsection{Sub Problem Solving}
For the final step $N$, $V(x_N)$ is known because $x_N$ is determined beforehand and $u_N=0$.
\begin{equation}
    V(x_N) = x_N^T P_N x_N \quad \text{where} \quad P_N = Q
\end{equation}
So, step $x_{N-1}$ could be solved and producing $V(t_{N-1})$ for the next step. Looping until the first step. In general, for each step,$u_t$ is the variable to be solved for \eqref{sub_problem}:
\begin{gather}
    V(x_t) = x_t^T P_t x_t\\
    \begin{bmatrix}
        x_t\\
        u_t
    \end{bmatrix} = X\\
    Q(x_t, u_t) = X^T\begin{bmatrix}Q, 0 \\ 0, R \end{bmatrix}X + x_{t+1}^T P_{t+1} x_{t+1}
\end{gather}
Considering system dynamics \eqref{system_dynamics}, $x_{t+1}$ is replaced by $x_t, u_t$:
\begin{gather}
    Q(x_t, u_t) = X^T\begin{bmatrix}Q, 0 \\ 0, R \end{bmatrix}X + (A x_t + B u_t)^T P_{t+1} (A x_t + B u_t)\\
    Q(x_t, u_t) = X^T\begin{bmatrix}Q+A^T P_{t+1} A, A^T P_{t+1} B \\ B^T P_{t+1} A, R + B^T P_{t+1} B \end{bmatrix}X\\
    Q(x_t, u_t) =X^T\begin{bmatrix}Q_{xx}, Q_{xu} \\ Q_{ux}, Q_{uu} \end{bmatrix}X\label{q_value_function}\\
    Q(x_t, u_t) = x_t^T Q_{xx} x_t + 2x_t^TQ_{xu} u_t + u_t^T Q_{uu} u_t
\end{gather}
Let gradient be zero:
\begin{equation}
    \frac{\partial{Q(x_t, u_t)}}{\partial{u_t}} = 2x_t^TQ_{xu} + 2Q_{uu} u_t = 0
\end{equation}
The optimal $\hat{u_t}$ is solved:
\begin{gather}
    \hat{u_t} = -Q_{uu}^{{-1}}Q_{ux} x_t = -K_t x_t\\
    Q_{xx} = Q+A^T P_{t+1} A\\
    Q_{xu} = A^T P_{t+1} B\\
    Q_{ux} = B^T P_{t+1} A\\
    Q_{uu} = R + B^T P_{t+1} B
\end{gather}
In the end, put the optimal control command into \eqref{q_value_function} to get $P_t$:
\begin{gather}
    V(x_t) = x_t^T P_t x_t = x_t^T(Q_{xx} - Q_{xu}Q_{uu}^{-1}Q_{ux})x_t\\
    P_t = Q_{xx} - Q_{xu}Q_{uu}^{-1}Q_{ux}
\end{gather}

\subsection{Forward Pass}
The procedure above is called "Backward Pass". after it is done, the only useful information is $K_t(t=0,1,...,N-1)$ which is a series of feedback gain. With this, we calculate the control command $u_t$ and update the states $x_t$ in the normal way, which is called "Forward Pass":\\
For $t=0,1,...,N-1$:\\
1. $u_t=-K_t x_t$\\
2. $x_{t+1} = A x_t + B u_t$\\

\end{document}
